% GENERATED FILE. Edit canonical lessons, then run npm run generate:books.
% canonical-source-hash: fnv1a64:246c5de7a9802073
% canonical-lessons: ML-C289-ankutti, ML-C289-ammayi, ML-C289-marumakal, ML-C289-ammayiyamma, ML-C289-ammayiyappan

\chapter{Boy, Aunt, Niece, Mother-in-law, Father-in-law}
\label{ch:ml-boy-aunt-niece-mother-in-law-father-in-law}
\hlchaptermodality{289}

\begin{chapteropening}
\textbf{By the end of this chapter:} \emph{I can say a boy, an aunt, a niece, a mother-in-law and a father-in-law.}

Complete the last lesson of chapter 289: I can say a boy, an aunt, a niece, a mother-in-law and a father-in-law.
\end{chapteropening}

\section[\texorpdfstring{\ml{ആൺകുട്ടി} (āṇkuṭṭi)}{ആൺകുട്ടി (āṇkuṭṭi)}]{\ml{ആൺകുട്ടി} (āṇkuṭṭi) — a boy}

\label{lesson:ML-C289-ankutti}



\begin{quote}
Before the new one: say the Malayalam for an amount, then the Malayalam for a form.
\end{quote}

\subsection*{You'll want to know: \ml{ആൺകുട്ടി}}
\textbf{\ml{ആൺകുട്ടി}} — \emph{āṇkuṭṭi} — \textquotedblleft{}a boy\textquotedblright{}.

\begin{grammarlens}[title={What you've built}]
One more word for this chapter.
\end{grammarlens}

\subsection*{Your turn}
\begin{itemize}
\raggedright
  \item \emph{Say it:} \emph{āṇkuṭṭi}
  \item \emph{Say it:} \emph{āṇkuṭṭi}, once more
  \item \emph{Say it:} say \emph{āṇkuṭṭi}
  \item \emph{Recall:} say the Malayalam for an amount, then the Malayalam for a form, then say \emph{āṇkuṭṭi} again
\end{itemize}

\begin{tcolorbox}[breakable,colback=teal!4,colframe=teal!35!black,title={Before you move on}]
What does \ml{ആൺകുട്ടി} mean? (\textquotedblleft{}A boy\textquotedblright{}.) Say it once more.
\end{tcolorbox}
\section[\texorpdfstring{\ml{അമ്മായി} (ammāyi)}{അമ്മായി (ammāyi)}]{\ml{അമ്മായി} (ammāyi) — an aunt (father's sister)}

\label{lesson:ML-C289-ammayi}



\begin{quote}
Before the new one: say the Malayalam for a form, then the Malayalam for a boy.
\end{quote}

\subsection*{You'll want to know: \ml{അമ്മായി}}
\textbf{\ml{അമ്മായി}} — \emph{ammāyi} — \textquotedblleft{}an aunt (father's sister)\textquotedblright{}.

\begin{grammarlens}[title={What you've built}]
One more word for this chapter.
\end{grammarlens}

\subsection*{Your turn}
\begin{itemize}
\raggedright
  \item \emph{Say it:} \emph{ammāyi}
  \item \emph{Say it:} \emph{ammāyi}, once more
  \item \emph{Say it:} say \emph{ammāyi}
  \item \emph{Recall:} say the Malayalam for a form, then the Malayalam for a boy, then say \emph{ammāyi} again
\end{itemize}

\begin{tcolorbox}[breakable,colback=teal!4,colframe=teal!35!black,title={Before you move on}]
What does \ml{അമ്മായി} mean? (\textquotedblleft{}An aunt (father's sister)\textquotedblright{}.) Say it once more.
\end{tcolorbox}
\section[\texorpdfstring{\ml{മരുമകൾ} (marumakaḷ)}{മരുമകൾ (marumakaḷ)}]{\ml{മരുമകൾ} (marumakaḷ) — a niece; a daughter-in-law}

\label{lesson:ML-C289-marumakal}



\begin{quote}
Before the new one: say the Malayalam for a boy, then the Malayalam for an aunt.
\end{quote}

\subsection*{You'll want to know: \ml{മരുമകൾ}}
\textbf{\ml{മരുമകൾ}} — \emph{marumakaḷ} — \textquotedblleft{}a niece; a daughter-in-law\textquotedblright{}.

\begin{grammarlens}[title={What you've built}]
One more word for this chapter.
\end{grammarlens}

\subsection*{Your turn}
\begin{itemize}
\raggedright
  \item \emph{Say it:} \emph{marumakaḷ}
  \item \emph{Say it:} \emph{marumakaḷ}, once more
  \item \emph{Say it:} say \emph{marumakaḷ}
  \item \emph{Recall:} say the Malayalam for a boy, then the Malayalam for an aunt, then say \emph{marumakaḷ} again
\end{itemize}

\begin{tcolorbox}[breakable,colback=teal!4,colframe=teal!35!black,title={Before you move on}]
What does \ml{മരുമകൾ} mean? (\textquotedblleft{}A niece; a daughter-in-law\textquotedblright{}.) Say it once more.
\end{tcolorbox}
\section[\texorpdfstring{\ml{അമ്മായിയമ്മ} (ammāyiyamma)}{അമ്മായിയമ്മ (ammāyiyamma)}]{\ml{അമ്മായിയമ്മ} (ammāyiyamma) — a mother-in-law}

\label{lesson:ML-C289-ammayiyamma}



\begin{quote}
Before the new one: say the Malayalam for an aunt, then the Malayalam for a niece.
\end{quote}

\subsection*{You'll want to know: \ml{അമ്മായിയമ്മ}}
\textbf{\ml{അമ്മായിയമ്മ}} — \emph{ammāyiyamma} — \textquotedblleft{}a mother-in-law\textquotedblright{}.

\begin{grammarlens}[title={What you've built}]
One more word for this chapter.
\end{grammarlens}

\subsection*{Your turn}
\begin{itemize}
\raggedright
  \item \emph{Say it:} \emph{ammāyiyamma}
  \item \emph{Say it:} \emph{ammāyiyamma}, once more
  \item \emph{Say it:} say \emph{ammāyiyamma}
  \item \emph{Recall:} say the Malayalam for an aunt, then the Malayalam for a niece, then say \emph{ammāyiyamma} again
\end{itemize}

\begin{tcolorbox}[breakable,colback=teal!4,colframe=teal!35!black,title={Before you move on}]
What does \ml{അമ്മായിയമ്മ} mean? (\textquotedblleft{}A mother-in-law\textquotedblright{}.) Say it once more.
\end{tcolorbox}
\section[\texorpdfstring{\ml{അമ്മായിയപ്പൻ} (ammāyiyappan)}{അമ്മായിയപ്പൻ (ammāyiyappan)}]{\ml{അമ്മായിയപ്പൻ} (ammāyiyappan) — a father-in-law}

\label{lesson:ML-C289-ammayiyappan}



\begin{quote}
Before the new one: say the Malayalam for a niece, then the Malayalam for a mother-in-law.
\end{quote}

\subsection*{You'll want to know: \ml{അമ്മായിയപ്പൻ}}
\textbf{\ml{അമ്മായിയപ്പൻ}} — \emph{ammāyiyappan} — \textquotedblleft{}a father-in-law\textquotedblright{}.

\begin{grammarlens}[title={What you've built}]
One more word for this chapter.
\end{grammarlens}

\subsection*{Your turn}
\begin{itemize}
\raggedright
  \item \emph{Say it:} \emph{ammāyiyappan}
  \item \emph{Say it:} \emph{ammāyiyappan}, once more
  \item \emph{Say it:} say \emph{ammāyiyappan}
  \item \emph{Recall:} say the Malayalam for a niece, then the Malayalam for a mother-in-law, then say \emph{ammāyiyappan} again
\end{itemize}

\begin{tcolorbox}[breakable,colback=teal!4,colframe=teal!35!black,title={Before you move on}]
What does \ml{അമ്മായിയപ്പൻ} mean? (\textquotedblleft{}A father-in-law\textquotedblright{}.) Say it once more.
\end{tcolorbox}
